\documentclass[12pt,letterpaper]{article}
\usepackage[T1]{fontenc}
\usepackage{newtxtext,newtxmath}
\usepackage[letterpaper,margin=1in]{geometry}
\usepackage{amsmath}
\usepackage{microtype}
\usepackage{booktabs,longtable,array,enumitem}
\usepackage[hidelinks]{hyperref}
\emergencystretch=2em
\title{Referee Report on\\\emph{On a metaplectic Jacquet--Langlands transfer\\for division algebras of prime degree}\\[.5em]\large Book Draft 2 (November 12, 2026)}
\author{Independent AI review by ChatGPT (GPT-6)}
\date{November 14, 2026}
\begin{document}
\maketitle
\noindent\textbf{Materials checked.} The GitHub \texttt{main} branch under
\url{https://github.com/wdjoyner/metaplectic-jl/tree/main/book},
including the master LaTeX file, Chapters 2--10, the appendix, bibliography,
and Draft 2 revision log; the attached \path{joyner-math-paper-style-guide.md}
(Draft 1, November 11, 2026); selected related published literature; and
\path{part-IV/scripts/level2_orbit_certificate.txt}.
The compiled book PDF was not independently inspected page by page, and the
full LaTeX source tree was not downloaded to a local build directory.
This is a selective mathematical re-refereeing, not a line-by-line proof
certification of the entire 216-page reported build.

\section{Recommendation}
\textbf{Verdict: promising and potentially original research monograph;
major mathematical fact-checking and editorial corrections remain before
it should be presented as a certified account.}
The work has a clear, substantial goal: a metaplectic analogue of the
Jacquet--Langlands correspondence for a central division algebra of prime
degree. The book correctly does not claim the full correspondence. Its
strongest specific contributions appear to be the ambient-cover and local
matching construction, the unitary division-side transform, the quadratic
toral positivity theorem, and the explicit wild cubic-cover character
calculations through level two. The global automorphic transfer retains
support and packet non-cancellation hypotheses; split-place and
archimedean comparison and wild integrality at arbitrary levels are open.

\section{Summary of results and proof status}
\begin{longtable}{@{}p{.12\linewidth}p{.56\linewidth}p{.23\linewidth}@{}}
\toprule Chapter & Principal result & Status presented\\\midrule\endhead
4 (I) & Realization of two covers inside one ambient cover; canonical
identification of good covered elliptic tori (A); matching away from the
center (B); packet-level trace identity and transfer (C). & Local results
proved using cited inputs; C conditional on $(\mathrm{Hyp}_D)$ and an
auxiliary supercuspidal place.\\
5 (II) & Canonical local $n$th-root section, germ transport near the center,
matching of arbitrary genuine test functions (D, E), and weakened
non-cancellation criterion (C$'$). & Matching proved using linear
harmonic-analysis inputs; global transfer conditional.\\
6 (III) & Unitary transform $L_D$ on central-character class functions,
tame irreducible-character integrality, a wild dimensional obstruction,
and local orbital-integral adjunction. & Tame results proved; general
wild integrality unresolved.\\
7 (IV) & Good-element/power-map calculations for tori, quadratic toral
permutation positivity, and cocycle-normalization and conditional global
comparison criteria. & Positivity presented as proved; global
comparison conditional.\\
8 (wild) & Cubic quaternion example over $F=\mathbb Q_3(\zeta_3)$:
Target A, dimension-27 genuine character from the trivial character;
Target B for representations trivial on $K^3$ (through level two),
including 54 degree-three intermediate characters. & Key Target A
components tagged (P$^c$); level-two proof uses hand arguments
and inherits the Target A computation.\\\bottomrule
\end{longtable}

\paragraph{Unitary character transform.}
For the specified cover of $H=D^\times$, the book constructs
$L_D:\mathcal H_\theta(H)\xrightarrow{\sim}
\mathcal H_{\widetilde\theta}(\widetilde H)$,
with normalization $b_D=|n|_F^{-(r-1)/2}$. If its irreducible-character
matrix is integral, unitarity and further sign arguments force a
permutation of irreducible characters. For $p\nmid n$ the book proves
integrality; for $p\mid n$ it is a genuinely new obstruction.

\paragraph{Toral theorem.}
For a quadratic extension $E/F$, $p$ odd, $n$ odd, and $\mu_n\subset F$,
the function $R_E(t)=\Delta(t)/\Delta(t^n)$, extended to scalars,
is shown to be a finite permutation character on $E^\times/F^\times$.
The proof is a radial valuation decomposition into regular characters of
finite quotients. This gives positive integral Fourier coefficients
without claiming extension of the toral modules to the entire
nonabelian cover.

\paragraph{Wild calculation.}
For $F=\mathbb Q_3(\zeta_3)$, $n=3$, $r=2$, the text finds
$b_D=3$ and $c_D=27$, exhibits $L_D\mathbf1$ as an irreducible
genuine character of dimension $27$, and proves the analogous
statement for representations trivial on $K^3$ (including norm twists).
The level-two case produces two irreducible constituents of dimension
$81$ on the relevant congruence subgroup.

\section{Originality assessment}
\textbf{Distinguish new proofs/results from new formulations of established
frameworks.} The ordinary Jacquet--Langlands correspondence is classical
\cite{JL70}; metaplectic correspondence for covered $\mathrm{GL}_r$ and
its test-function theory precede this book \cite{FK86}. Joyner's 1988
paper already studied covers attached to division algebras and metaplectic
lifting \cite{Jo88}. Zou corrected and developed the local
Flicker--Kazhdan correspondence \cite{Zo23}; Kaplan--Lapid--Zou
classified irreducibles for covers of $\mathrm{GL}_r$ \cite{KLZ23}.
Thus neither covering groups, the basic power map, nor the general
trace-formula strategy are themselves original here.

\begin{description}[style=nextline,leftmargin=0pt]
\item[High potential originality.] The concrete division-side wild
integrality result, its finite-design/Clifford orbit machinery, and the
uniform quadratic toral permutation-character theorem appear distinctive
relative to the works consulted. Priority requires a more exhaustive
literature search and independent proof verification.
\item[Moderate to high potential originality.] The ambient common-cover
identification and two-sided matching for arbitrary genuine test
functions near the center provide new work in this specific division-algebra
setting, though their ingredients are established tools.
\item[Moderate originality.] The unitary $L_D$ transform and tame
integrality theorem may be mathematically useful new formulations
and proofs, but their relationship to the 1988 division-side
character relations should be sharpened to delimit the new claims.
\item[More limited originality.] The conditional global transfer and
cocycle normalization are substantial synthesis and adaptation,
but the trace-formula and representation-separation techniques are
classical; a fully unconditional transfer is not yet established.
\end{description}

\noindent\emph{Originality verdict.} The work should be described as a
partial program with several candidate new theorems rather than as
a completed new correspondence. Search results did not locate an earlier
identical result for the explicit wild computations, but absence of a
search result does not establish priority.

\section{Verification performed and limits}
Following the attached guide, the levels are: 1 = independent derivation;
2 = checked against an actually consulted source; 3 = background knowledge
only. Proof status labels in the book are distinct from these levels.

\paragraph{Level 1: arithmetic.} In $\mathbb Z[\zeta_3]$,
$4-\zeta_3=(1-\zeta_3)(3+\zeta_3)$, and the three norms are $21,3,7$,
respectively. This independently supports the stated factorization of
the ramified primes. Since $|3|_{\mathbb Q_3(\zeta_3)}=3^{-2}$,
$b_D=3$ and $c_D=27$ follow directly.

\paragraph{Level 1: toral cubic example.} With residue field size $3$,
the valuation formula in Chapter 8 gives
$R_{\rm ur}=1+2\operatorname{Reg}(C_4)$ and
$R_{\rm ram}=1+\operatorname{Reg}(C_2)+\operatorname{Reg}(C_6)$.
Both have value $9$ at the identity, as predicted by $|3|_F^{-1}$.
This validates these numerical instances, not the full general theorem.

\paragraph{Level 1: independent certificate consistency.} I obtained
the committed 81,944-character level-two orbit certificate via GitHub
and checked its tables independently, using fresh group-table and
action-table assertions rather than accepting the manuscript's report.
The alleged group of order $27$ is associative with identity; both
listed permutation actions satisfy the group-action property and have
no nonidentity fixed point. The 54 degree-three characters form two
orbits of size 27, and the 243 linear characters form nine orbits of
size 27. This certifies \emph{internal consistency of the tables},
not their derivation from the Hilbert-symbol cocycle.

\paragraph{Level 2: external literature.} Publication metadata and
abstracts for \cite{FK86,Jo88,Zo23,KLZ23} were consulted. Their entire
proofs and the theorem-specific dependencies in the book were not
independently cross-checked. No Level 1 check of the global trace-formula
proof was completed. Load-bearing design and Hilbert-symbol scripts
were inspected for provenance but not executed in full. Accordingly,
(P$^c$) statements remain conditional on the certificate and computations
being correct as claimed.

\section{Compliance with the attached style guide}
\textbf{Substantially aligned in structure, but not fully compliant in
verification and exposition.} The main source uses
\verb|\documentclass[12pt,letterpaper,oneside]{book}|, Times text and
math via \texttt{newtx} with \texttt{mathptmx} fallback, and one-inch
margins. Keeping the existing \texttt{book} class follows the guide's
[DECIDE] default. The required introduction, actual literature review,
background and notation, body, summary, glossary, acknowledgments,
references and revision log are present. The email and AI roles are
correctly stated, and the acknowledgments explicitly reject the notion
that AI review certifies the work.

A static scan of the master and twelve included text files found
495 distinct \verb|\label| values, no unresolved \verb|\ref|/
\verb|\eqref|/\verb|\cref| targets in the source tree, and no missing
\verb|\cite| keys among the 106 \verb|\bibitem| entries. This is a
source-level consistency audit, not a proof of successful compilation.
The Draft 2 revision log reports a successful 216-page build with no
undefined references or overfull boxes under the available Times
fallback. I did not independently rebuild the book or audit the
rendered pages.

The number of worked examples is modest relative to the number of
technical sections: Chapters 4--8 contain 23 explicit
\verb|example| environments in total, while many subsections have no
standalone worked example. Some sections need no example, as the author's
November 12 ruling records, but further examples in the wild and global
sections would materially help readers. Long paragraphs and semicolon-heavy
sentences remain, especially in Chapters 4, 5 and 8. The guide prefers
short sentences. Several displayed named statuses use parentheses
(P), (E), (O), (C), rather than the square-bracket notation in the
guide; this is an intentional but undocumented choice relative to the
attached version, and needs a decision rather than being treated as a
mathematical defect. The book's chapter-and-section theorem numbering
(4.4.5, etc.) is also intentional and explicitly authorized in the
revision log despite the paper guide's simpler example.

The revision log explicitly says citations were not rechecked and scripts
were not rerun during the Draft 2 pass. This means the full verification
requirements of style-guide Sections 4.3--4.5 remain outstanding, even
though the manuscript correctly records its limitations. In a revised
edition, provide an audit table identifying each substantial imported
theorem, exact theorem/page citation, and independent check; attach a
machine-readable log for every load-bearing computation, with source
commit and dependency versions.

\section{Relevant earlier literature}
\begin{thebibliography}{KLZ23}
\bibitem[JL70]{JL70} H. Jacquet and R. P. Langlands,
\emph{Automorphic Forms on GL(2)}, Springer Lecture Notes in Mathematics 114 (1970).
\bibitem[FK86]{FK86} Y. Z. Flicker and D. A. Kazhdan,
\emph{Metaplectic correspondence}, Publ. Math. IH\'ES \textbf{64} (1986), 53--110.
\bibitem[Jo88]{Jo88} D. Joyner,
\emph{On the metaplectic analog of Kazhdan's ``endoscopic'' lifting},
Israel J. Math. \textbf{61} (1988), 113--154.
\bibitem[Zo23]{Zo23} J. Zou,
\emph{Local metaplectic correspondence and applications},
Math. Z. \textbf{305} (2023), article 43.
\bibitem[KLZ23]{KLZ23} E. Kaplan, E. Lapid and J. Zou,
\emph{Classification of irreducible representations of metaplectic
covers of the general linear group over a non-archimedean local field},
Represent. Theory \textbf{27} (2023), 1041--1087.
\end{thebibliography}

\section{Errors and required corrections (listed last)}
\begin{enumerate}[leftmargin=*,label=\textbf{E\arabic*.}]
\item \textbf{False statement, verified independently (Level 1): local
division algebras are not unique in degree $m>2$.}
\texttt{book/chapters/ch3\_notation.tex}, about lines 127--133,
says that for each degree $m$ there is exactly one central division
algebra over a nonarchimedean local field, and that invariants
$1/m,-1/m,k/m$ designate the same algebra. This is false.
$\operatorname{Br}(F)\simeq\mathbb Q/\mathbb Z$ distinguishes them:
the number of isomorphism classes with degree $m$ is $\varphi(m)$;
for $m=3$, invariants $1/3$ and $2/3$ give distinct division
algebras (opposites). Replace the passage by a correct classification
and adapt the cyclic construction to Frobenius power $\sigma^a$.
Chapter 8, lines 156--170, already states this correctly and supplies
a consistency check within the book. This error lies in the background,
not a demonstrated counterexample to an original theorem.

\item \textbf{False explanatory assertion (Level 1): cocycle existence
versus scalar centrality.} \texttt{book/mjl\_book\_draft.tex},
Introduction, ``The two covers'', says $n\mid r-1+2rc_v$ is exactly
the condition making the local Kazhdan--Patterson cocycle well defined.
But $\sigma_F(g,h)(\det g,\det h)_{n,F}^{c_v}$ is a 2-cocycle
for every integer $c_v$ because it is a product of 2-cocycles.
The displayed divisibility condition instead makes the full preimage
of scalar matrices central; this is explicitly the use made of it
in Chapter 4, Lemma \texttt{I:lem:center}, and Chapter 7's diagonal
commutator audit. Replace ``well defined'' by ``has central scalar
preimage'' (under the specified conventions).

\item \textbf{Overclaim (fixable misstatement, Level 2): wild obstruction
not completely resolved.} The Introduction's chapter overview and
running-example discussion describe Chapter 8 as resolving the
Chapter 6 wild obstruction ``completely'' in the example. But the
same book proves Target A and Target B only for representations
trivial on $K^3$ and explicitly leaves higher congruence levels
open (Chapter 8, Future directions). Replace with ``resolves the
trivial-character case and all cases through level two''.

\item \textbf{Unverified dependencies (fixable verification gap, not a
proved mathematical error).} Full numerical and cocycle audit for
Target A is not repeated here; the whole 216-page book has not
been independently recompiled, and its imported theorem citations
have not all been checked theorem by theorem. Retain the (P$^c$)
labels until an independent rerun of the exact-design and symbol
scripts confirms their outputs, and record precisely what was tested.
\end{enumerate}
\end{document}
